\chapter{测量偏倚与非因果图}

% 本章多张图共用“处理—结局各自一条测量链”的骨架
\newcommand{\measbase}{%
  \node[wvtx] (A)  at (0,0)       {$A$};
  \node[wvtx] (Ap) at (0,1.15)    {$A^{*}$};
  \node[wlat] (UA) at (0,2.3)     {$U_A$};
  \node[wvtx] (Y)  at (3.3,0)     {$Y$};
  \node[wvtx] (Yp) at (3.3,1.15)  {$Y^{*}$};
  \node[wlat] (UY) at (3.3,2.3)   {$U_Y$};
  \draw[ed] (UA) -- (Ap);
  \draw[ed] (Ap) -- (A);
  \draw[ed] (UY) -- (Yp);
  \draw[ed] (Yp) -- (Y);
  \draw[ed] (A) -- (Y);
}

% 本章节点标签含上下标，用本章专用节点样式（不改动共享的 vtx/lat）
\tikzset{
  wvtx/.style={circle,draw,inner sep=2.4pt,minimum size=7.6mm,font=\small},
  wlat/.style={circle,draw,dashed,inner sep=2.4pt,minimum size=7.6mm,font=\small},
}

设想一位研究者做了一项随机化实验，想回答这样一个因果问题：“一个人抬头望天，会不会让别的行人也抬头望天？”她发现，自己抬头与其他行人抬头之间只有微弱的关联。这一微弱的关联反映的是微弱的因果效应吗？按随机化实验的定义，本研究中不应存在混杂偏倚。此外也不应存在选择偏倚，因为所有行人的反应——无论抬头与否——都被记录了下来。然而还有另一个问题：负责记录行人反应的那位合作者出了很多错。具体地说，行人抬头的场合有一半被她漏掉了，并记成了“没有抬头”。于是，即使处理（研究者抬头）对结局（他人抬头）真有很强的效应，结局的错分也会把处理与（测量有误的）结局之间的关联稀释掉。

当研究数据的测量过程使处理与结局之间的关联被削弱或增强时，我们说存在\emph{测量偏倚}（measurement bias）。由于测量误差在任何研究设计中都可能出现——随机化实验与观察性研究皆然——解读效应估计值时，永远要考虑测量偏倚。本章描述由测量误差引起的偏倚。

\section{测量误差}

此前各章隐含地作了一个不现实的假定：所有变量都被完美测量。考虑一项观察性研究，目的是估计降胆固醇药物 $A$ 对肝病风险 $Y$ 的效应。我们往往预期处理 $A$ 会被不完美地测量。例如，若用药信息得自病历摘录，摘录者可能抄错数据，医生可能忘记写下曾给病人开过该药，病人也可能没有按处方服药。于是，分析数据集中的处理变量不会是真实的用药，而是测量到的用药。我们把测量到的处理记作 $A^{*}$（读作 A-star），对某个个体而言，它未必等于真实处理 $A$。心理学文献有时把 $A$ 称为“构念”（construct），把 $A^{*}$ 称为“测量”或“指标”。观察性学科面临的挑战，正是用观测到的测量（例如病历中的他汀使用信息）去推断未被观测的构念（例如降胆固醇药物的使用）。

图 9.1 的因果图描绘了变量 $A$、$A^{*}$ 与 $Y$。为简单起见，我们选了一个对 $A$ 作用于 $Y$ 的因果效应既无混杂、也无选择偏倚的情形。真实处理 $A$ 同时影响结局 $Y$ 与测量到的处理 $A^{*}$。图中还包含节点 $U_A$，代表除 $A$ 之外决定 $A^{*}$ 取值的一切因素。我们把一个个体的测量值 $A^{*}$ 与真实值 $A$ 之差，称为该个体 $A$ 的\emph{测量误差}。测量误差的大小与方向由 $U_A$ 中的因素决定。

\mnote{离散变量的测量误差称为\emph{错分}（misclassification）。}

注意，把节点 $U_A$ 画进因果图并非严格必需，因为 $U_A$ 既不是图中其他变量共享的原因，也不是被条件化的变量。我们之所以画它，是为了显式地表示造成测量误差的因素，并与下文将要讨论的因果图直接对照。

\begin{figure}[htbp]
\centering
\begin{tikzpicture}[causal,scale=0.95,every node/.style={transform shape}]
  \node[wlat] (UA) at (0,2.4)  {$U_A$};
  \node[wvtx] (Ap) at (0,1.2)  {$A^{*}$};
  \node[wvtx] (A)  at (0,0)    {$A$};
  \node[wvtx] (Y)  at (2.9,0)  {$Y$};
  \draw[ed] (UA) -- (Ap);
  \draw[ed] (Ap) -- (A);
  \draw[ed] (A) -- (Y);
\end{tikzpicture}
\caption{处理 $A$ 的测量误差：真实处理 $A$ 决定测量到的处理 $A^{*}$，$U_A$ 表示造成测量误差的因素。}
\label{fig:9.1}
\end{figure}

\renewcommand{\TPname}{技术注记（Technical Point）}
\begin{techpoint}[测量误差的独立性与非差异性]
对每个个体，我们把 $A$ 的测量误差定义为差值 $e_A=A^{*}-A$，把 $Y$ 的测量误差定义为差值 $e_Y=Y^{*}-Y$。

以 $f(\cdot)$ 表示概率密度函数（pdf）。处理 $A$ 的测量误差 $e_A$ 与结局 $Y$ 的测量误差 $e_Y$ \emph{独立}，是指二者的联合 pdf 等于各自边缘 pdf 之积，即 $f(e_Y,e_A)=f(e_Y)f(e_A)$。处理 $A$ 的测量误差 $e_A$ \emph{非差异}（nondifferential），是指其 pdf 与结局 $Y$ 无关，即 $f(e_A\given Y)=f(e_A)$。类似地，结局 $Y$ 的测量误差 $e_Y$ 非差异，是指其 pdf 与处理 $A$ 无关，即 $f(e_Y\given A)=f(e_Y)$。
\end{techpoint}
\renewcommand{\TPname}{技术注记}

除处理 $A$ 之外，结局 $Y$ 也可能被测量出错。图 9.2 的因果图包含测量到的结局 $Y^{*}$，以及造成 $Y$ 之测量误差的因素 $U_Y$。图 9.2 刻画了实践中的常见情形：人们想计算处理 $A$ 对结局 $Y$ 的平均因果效应，但这两个变量 $A$ 与 $Y$ 从未被正确测量，也无法被正确测量；研究者手上只有测量有误的版本 $A^{*}$ 与 $Y^{*}$，却要识别 $A$ 对 $Y$ 的因果效应。

\begin{figure}[htbp]
\centering
\begin{tikzpicture}[causal,scale=0.95,every node/.style={transform shape}]
  \measbase
\end{tikzpicture}
\caption{处理与结局均有测量误差，且两项测量误差彼此独立。}
\label{fig:9.2}
\end{figure}

图 9.2 还表示了一个对处理 $A$ 作用于结局 $Y$ 的因果效应既无混杂、也无选择偏倚的情形。按前几章的推理，在这种情形下关联即因果。我们可以计算任何 $A$–$Y$ 关联度量，并赋予其因果解释，即把它当作 $A$ 对 $Y$ 的效应。例如，关联风险比 $\Pr[Y=1\given A=1]$ / $\Pr[Y=1\given A=0]$ 等于因果风险比 $\Pr[Y^{a=1}=1]$ / $\Pr[Y^{a=0}=1]$。此前各章隐含的假定——现在我们把它明确说出来——是我们拥有 $A$ 与 $Y$ 被完美测量的数据。

现在考虑更现实的情形：只有测量有误的处理与结局，即 $A^{*}$ 与 $Y^{*}$。这时无法保证 $A^{*}$ 与 $Y^{*}$ 之间的关联度量等于 $A$ 对 $Y$ 的因果效应度量。关联风险比 $\Pr[Y^{*}=1\given A^{*}=1]$ / $\Pr[Y^{*}=1\given A^{*}=0]$ 一般不等于因果风险比 $\Pr[Y^{a=1}=1]$ / $\Pr[Y^{a=0}=1]$。我们说存在\emph{测量偏倚}或\emph{信息偏倚}（information bias）。存在测量偏倚时，可交换性、正值性与一致性这三条可识别性条件，不足以为我们算出处理 $A$ 对结局 $Y$ 的因果效应。

\section{测量误差的结构}

混杂的因果结构可以概括为：处理与结局存在共同原因。选择偏倚的因果结构则可以概括为：对处理与结局的共同效应（或其原因）作了条件化。测量偏倚在存在测量误差时出现，但测量误差没有单一的结构可以概括。本节按两种性质——独立性与非差异性——对测量误差的结构进行分类，这两种性质在下文描述（形式定义见技术注记 9.1）。

\begin{figure}[htbp]
\centering
\begin{tikzpicture}[causal,scale=0.95,every node/.style={transform shape}]
  \measbase
  \node[wlat] (UAY) at (1.65,3.2) {$U_{AY}$};
  \draw[ed] (UAY) -- (UA);
  \draw[ed] (UAY) -- (UY);
\end{tikzpicture}
\caption{处理与结局的测量误差相依：$U_{AY}$ 同时影响两项测量。}
\label{fig:9.3}
\end{figure}

图 9.2 的因果图描绘了处理 $A$ 与结局 $Y$ 各自的测量误差 $U_A$ 与 $U_Y$。按 d-分离的规则，测量误差 $U_A$ 与 $U_Y$ 相互独立，因为二者之间的路径被对撞节点（$A^{*}$ 或 $Y^{*}$）阻断。若例如用药 $A$ 与肝毒性 $Y$ 的信息都取自电子病历，而数据录入错误是随意发生的，就可以预期出现独立的测量误差。然而在另一些情形下，暴露与结局的测量误差可能相依，如图 9.3 所示。例如，若信息是事后通过电话访谈获得的，而个体回忆自己病史的能力（$U_{AY}$）同时影响处理 $A$ 与结局 $Y$ 的测量，就会出现相依的测量误差。

图 9.2 与图 9.3 表示的情形是：造成处理之测量误差的因素 $U_A$ 与结局的真实取值 $Y$ 独立，造成结局之测量误差的因素 $U_Y$ 与处理的真实取值 $A$ 独立。这时我们说，处理的测量误差相对于结局是\emph{非差异的}，结局的测量误差相对于处理也是非差异的。图 9.4 的因果图给出了独立但\emph{差异}（differential）的测量误差之一例：结局的真实取值影响了处理的测量（即存在从 $Y$ 指向 $U_A$ 的箭头）。下面举几个处理之差异性测量误差的例子。

\begin{figure}[htbp]
\centering
\begin{tikzpicture}[causal,scale=0.95,every node/.style={transform shape}]
  \measbase
  \draw[ed] (Y) -- (UA);
\end{tikzpicture}
\caption{独立但差异的测量误差：结局的真实取值影响处理的测量。}
\label{fig:9.4}
\end{figure}

设结局 $Y$ 不是肝毒性而是痴呆，用药 $A$ 通过访谈研究对象来获知。由于痴呆会影响回忆 $A$ 的能力，可以预期存在从 $Y$ 指向 $U_A$ 的箭头。同样，在一项计算孕期饮酒 $A$ 对出生缺陷 $Y$ 之效应的研究中，若饮酒量是在分娩后凭回忆获知的，也可以预期存在从 $Y$ 指向 $U_A$ 的箭头——因为回忆可能受妊娠结局的影响。这两个例子中由此产生的测量偏倚，常被称为\emph{回忆偏倚}（recall bias）。若用药物 $A^{*}$ 的血药浓度代替实际用药 $A$，而血药浓度是在肝毒性 $Y$ 出现之后测得的，也可能产生结构相同的偏倚——因为肝毒性会影响所测得的血药浓度。由此产生的测量偏倚常被称为\emph{反向因果偏倚}（reverse causation bias）。

\begin{figure}[htbp]
\centering
\begin{tikzpicture}[causal,scale=0.95,every node/.style={transform shape}]
  \measbase
  \draw[ed] (A) -- (UY);
\end{tikzpicture}
\caption{独立但差异的测量误差：处理的真实取值影响结局的测量。}
\label{fig:9.5}
\end{figure}

图 9.5 的因果图给出了独立但差异的测量误差之一例：处理的真实取值影响了结局的测量（即存在从 $A$ 指向 $U_Y$ 的箭头）。若医生怀疑用药 $A$ 引起肝毒性 $Y$，因而对用药病人比其他人监测得更密切，就会出现结局的差异测量误差。图 9.6 与图 9.7 描绘的测量误差既相依又差异，可能由上述情形的组合造成。

\begin{figure}[htbp]
\centering
\begin{tikzpicture}[causal,scale=0.95,every node/.style={transform shape}]
  \measbase
  \node[wlat] (UAY) at (1.65,3.2) {$U_{AY}$};
  \draw[ed] (UAY) -- (UA);
  \draw[ed] (UAY) -- (UY);
  \draw[ed] (Y) -- (UA);
\end{tikzpicture}
\caption{既相依又差异的测量误差：$U_{AY}$ 使两项测量相依，结局又影响处理的测量。}
\label{fig:9.6}
\end{figure}

\begin{figure}[htbp]
\centering
\begin{tikzpicture}[causal,scale=0.95,every node/.style={transform shape}]
  \measbase
  \node[wlat] (UAY) at (1.65,3.2) {$U_{AY}$};
  \draw[ed] (UAY) -- (UA);
  \draw[ed] (UAY) -- (UY);
  \draw[ed] (A) -- (UY);
\end{tikzpicture}
\caption{既相依又差异的测量误差：$U_{AY}$ 使两项测量相依，处理又影响结局的测量。}
\label{fig:9.7}
\end{figure}

总之，我们讨论了四类测量误差：独立非差异（图 9.2）、相依非差异（图 9.3）、独立差异（图 9.4 与图 9.5）、相依差异（图 9.6 与图 9.7）。测量误差的具体结构，决定了可用哪些方法加以校正。例如，当测量误差为独立非差异时，关于测量误差校正方法的文献非常丰富。一般说来，测量误差的校正方法依赖建模假定与验证样本（validation samples）的结合；所谓验证样本，是指其中关键变量测得几乎没有误差的那部分数据。测量误差校正方法的细节超出本书范围。我们的目的只是强调：测量变量这一行为（如同抽取个体这一行为）可能引入偏倚（关于其强度与方向，见精要 9.1）。贴近现实的因果图需要同时表示由混杂、选择与测量引起的偏倚。对抗测量有误所致偏倚的最佳办法，显然是改进分析中所用变量的测量程序。

\renewcommand{\FPname}{精要（Fine Point）}
\begin{finepoint}[测量偏倚的强度与方向]
一般说来，测量误差会导致偏倚。一个值得注意的例外是 $A$ 与 $Y$ 无关联、且测量误差为独立非差异的情形：若图 9.2 中从 $A$ 到 $Y$ 的箭头不存在，则 $A$–$Y$ 关联与 $A^{*}$–$Y^{*}$ 关联都将为零。在其他一切情形下，测量偏倚都可能使 $A^{*}$–$Y^{*}$ 关联比 $A$–$Y$ 关联离零更远或更近。更糟的是，即使处于图 9.2 那种独立非差异的测量误差结构之下，非极端的测量偏倚也可能使 $A^{*}$–$Y^{*}$ 与 $A$–$Y$ 的趋势相反——对非二值的有序处理与连续处理即是如此。这种独立非差异测量误差之下的趋势反转，出现在给定 $A$ 时 $A^{*}$ 的条件均值是 $A$ 的非单调函数之时。细节见 Dosemeci, Wacholder, and Lubin (1990) 与 Weinberg, Umbach, and Greenland (1994)。VanderWeele and Hernán (2009) 用带符号的因果图（signed causal diagrams）描述了一个更一般的框架。

测量偏倚的大小取决于测量误差的大小。也就是说，测量偏倚一般随 $U_A\to A^{*}$ 与 $U_Y\to Y^{*}$ 这两条箭头的强度而增大。因果图不编码定量信息，因此不能用它来描述偏倚的大小。
\end{finepoint}
\renewcommand{\FPname}{精要}

\section{测量有误的混杂因子与对撞结构}

除处理 $A$ 与结局 $Y$ 之外，混杂因子 $L$ 也可能被测量出错。即使处理与结局都被完美测量，混杂因子测量有误也可能导致偏倚。

为看清这一点，考虑图 9.8 的因果图，其中包含变量用药 $A$、肝病 $Y$ 与肝炎病史 $L$。有既往肝炎 $L$ 的人更不容易被处方药物 $A$，也更可能发生肝病 $Y$。如第 7 章所述，处理 $A$ 对结局 $Y$ 的效应存在混杂，因为存在一条开放的后门路径 $A\leftarrow L\rightarrow Y$；但给定 $L$ 后不存在未测混杂，因为后门路径 $A\leftarrow L\rightarrow Y$ 可以通过对 $L$ 条件化而阻断。也就是说，给定混杂因子 $L$，受治疗者与未受治疗者是可交换的，可以应用逆概率加权或标准化来计算 $A$ 对 $Y$ 的平均因果效应。基于 $L$、$Y$ 与 $A$ 的标准化（或逆概率加权）风险比，将等于因果风险比 $\Pr[Y^{a=1}=1]$ / $\Pr[Y^{a=0}=1]$。

\begin{figure}[htbp]
\centering
\begin{tikzpicture}[causal,scale=0.95,every node/.style={transform shape}]
  \node[wvtx] (L)  at (0,0)     {$L$};
  \node[wvtx] (Lp) at (0,1.35)  {$L^{*}$};
  \node[wvtx] (A)  at (1.9,0)   {$A$};
  \node[wvtx] (Y)  at (3.7,0)   {$Y$};
  \draw[ed] (L) -- (Lp);
  \draw[ed] (L) -- (A);
  \draw[ed] (A) -- (Y);
  \draw[ed] (L.east) .. controls (1.0,1.45) and (2.6,1.45) .. (Y.west);
\end{tikzpicture}
\caption{混杂因子 $L$ 及其测量值 $L^{*}$。}
\label{fig:9.8}
\end{figure}

上述推理中隐含的假定同样是：混杂因子 $L$ 被完美测量。设研究者无法获得研究对象的病历。为了确认既往肝炎诊断，研究者只能通过问卷询问研究对象。由于并非所有对象都能准确回忆自己的病史——有人不愿让人知道，有人记性不好，有人只是答错了问卷——混杂因子 $L$ 被测量出错。注意，图 9.8 并未显式表示造成 $L$ 之测量误差的因素 $U_L$，因为这一误差的具体结构与本处讨论无关。

研究者手上只有测量有误的变量 $L^{*}$，而没有 $L$。不幸的是，后门路径 $A\leftarrow L\rightarrow Y$ 一般不能通过对 $L^{*}$ 条件化而阻断。基于 $L^{*}$、$Y$ 与 $A$ 的标准化（或逆概率加权）风险比，一般不等于因果风险比 $\Pr[Y^{a=1}=1]$ / $\Pr[Y^{a=0}=1]$。这时我们说存在测量偏倚或信息偏倚。图 9.9 的因果图给出了 $A$ 对 $Y$ 之因果效应存在混杂、而 $L$ 并非 $A$ 与 $Y$ 共有原因的一例。这里 $L$ 的测量有误同样导致测量偏倚，因为后门路径 $A\leftarrow L\leftarrow U\rightarrow Y$ 一般不能通过对 $L^{*}$ 条件化而阻断。

\begin{figure}[htbp]
\centering
\begin{tikzpicture}[causal,scale=0.95,every node/.style={transform shape}]
  \node[wlat] (U)  at (0,-1.7)  {$U$};
  \node[wvtx] (L)  at (0,0)     {$L$};
  \node[wvtx] (Lp) at (0,1.35)  {$L^{*}$};
  \node[wvtx] (A)  at (1.9,0)   {$A$};
  \node[wvtx] (Y)  at (3.7,0)   {$Y$};
  \draw[ed] (L) -- (Lp);
  \draw[ed] (L) -- (A);
  \draw[ed] (A) -- (Y);
  \draw[ed] (U) -- (L);
  \draw[ed] (U) -- (Y);
\end{tikzpicture}
\caption{$L$ 并非 $A$ 与 $Y$ 的共有原因时，$L$ 的测量误差仍导致测量偏倚。}
\label{fig:9.9}
\end{figure}

另一种看法是，把图 9.8 与图 9.9 中因混杂因子测量有误所致的偏倚，视为一种未测混杂，而不是一种测量偏倚。事实上，图 9.8 的因果图等价于图 7.8。可以把 $L$ 看作未测变量，把 $L^{*}$ 看作替代混杂因子（surrogate confounder，见精要 7.2）。究竟用哪种说法——未测混杂，还是混杂因子测量有误所致偏倚——对实践目的无关紧要。不过在有些情形下，使用测量有误的变量已足以校正混杂，例子见精要 9.2。

混杂因子测量有误还可能造成表观的效应修饰。举例来说，设所有报告有既往肝炎诊断（$L^{*}=1$）的研究对象，以及报告无既往肝炎诊断（$L^{*}=0$）者中的一半，实际上确有既往肝炎诊断（$L=1$）。也就是说，混杂因子的真值与测量值在层 $L^{*}=1$ 中一致，在层 $L^{*}=0$ 中不一致。再设处理 $A$ 对任何个体的肝病 $Y$ 都无效应，即严格零假设成立。当研究者把分析限制在层 $L^{*}=1$ 时，不存在由 $L$ 造成的混杂，因为纳入分析的对象都具有相同的 $L$ 取值（即 $L=1$）。因此在层 $L^{*}=1$ 中，他们找不到 $A$ 与 $Y$ 之间的关联。然而，当研究者把分析限制在层 $L^{*}=0$ 时，就存在由 $L$ 造成的混杂，因为层 $L^{*}=0$ 中既有 $L=1$ 的人也有 $L=0$ 的人。于是，由于 $L$ 的混杂未受控制，研究者会发现 $A$ 与 $Y$ 之间存在非零关联。如果研究者并不知道混杂因子在层 $L^{*}=0$ 中测量有误、而在层 $L^{*}=1$ 中没有测量误差，他们就可能天真地得出结论：层 $L^{*}=0$ 的关联度量与层 $L^{*}=1$ 的关联度量都可以解释为效应度量。由于这两个关联度量不同，研究者会说 $L^{*}$ 是 $A$ 对 $Y$ 之效应的修饰因子，尽管真实的混杂因子 $L$ 并不存在效应修饰。

最后，对撞结构 $C$ 也可能被测量出错。这一情形见图 9.10 的因果图，它等价于图 8.2。在图 9.10 之下关心 $A$ 对 $Y$ 的效应时，对测量有误的对撞节点 $C^{*}$ 条件化一般会引入选择偏倚，因为 $C^{*}$ 是对撞节点的后代，从而是处理 $A$ 与结局 $Y$ 的共同效应。

\begin{figure}[htbp]
\centering
\begin{tikzpicture}[causal,scale=0.95,every node/.style={transform shape}]
  \node[wvtx] (A)  at (0,0)    {$A$};
  \node[wvtx] (Y)  at (1.35,0) {$Y$};
  \node[wvtx] (C)  at (2.7,0)  {$C$};
  \node[draw,rectangle,inner sep=1.5pt,minimum size=7mm,font=\small] (Cp) at (4.1,0) {$C^{*}$};
  \draw[ed] (A) -- (Y);
  \draw[ed] (Y) -- (C);
  \draw[ed] (C) -- (Cp);
  \draw[ed] (A) .. controls (0.9,1.15) and (1.9,1.15) .. (C);
\end{tikzpicture}
\caption{对撞结构 $C$ 被测量出错：$C^{*}$ 为 $C$ 的测量值。}
\label{fig:9.10}
\end{figure}

\begin{finepoint}[测量有误的混杂因子何时不成问题]
在许多医学应用中，混杂因子的测量误差不会引入任何偏倚。设高血压 $L$ 同时影响接受降压治疗 $A$ 的概率与发生卒中 $Y$ 的概率。然而医生与病人作治疗决定时依据的并不是真实血压 $L$，而是记录在诊室病历里的血压测量值 $L^{*}$。也就是说，$L^{*}$ 是决策者所能获得的关于血压的唯一信息。图 9.11（结构上等价于图 7.2）表示这一情形。可能测量有误的 $L^{*}$ 完全中介了 $L$ 对 $A$ 的效应，因为 $L$ 中未被 $L^{*}$ 捕捉到的任何成分都无人知晓，因而无法影响给药决定。由此可知，$A$ 与 $Y$ 之间的后门路径既可以通过对真实 $L$ 条件化而阻断，也可以通过对测量到的 $L^{*}$ 条件化而阻断。因此，研究者拿到的是真实 $L$ 还是测量到的 $L^{*}$，都无关紧要；两个变量中的任何一个都足以校正混杂。

\begin{center}
\begin{tikzpicture}[causal,scale=0.95,every node/.style={transform shape}]
  \node[wvtx] (L)  at (0,-1.7) {$L$};
  \node[wvtx] (Lp) at (0,0)    {$L^{*}$};
  \node[wvtx] (A)  at (2.0,0)  {$A$};
  \node[wvtx] (Y)  at (3.9,0)  {$Y$};
  \draw[ed] (L) -- (Lp);
  \draw[ed] (Lp) -- (A);
  \draw[ed] (A) -- (Y);
  \draw[ed] (L) -- (Y);
\end{tikzpicture}
\captionof{figure}{$L^{*}$ 完全中介 $L$ 对 $A$ 的效应：校正 $L$ 或校正 $L^{*}$ 均可。}
\label{fig:9.11}
\end{center}

图 9.12 给出了一个更极端的例子。在该因果图之下，拥有真实 $L$ 的数据不足以校正混杂，而拥有测量到的 $L^{*}$ 的数据却足以校正混杂。要点在于：只要数据中的信息量达到决策者作决定时所拥有的信息量，效应就是可识别的——无论这些信息来自完美测量的变量，还是来自测量有误的变量。

\begin{center}
\begin{tikzpicture}[causal,scale=0.92,every node/.style={transform shape}]
  \node[wlat] (U1) at (0,2.5)  {$U_1$};
  \node[wlat] (U2) at (0,-1.7) {$U_2$};
  \node[wvtx] (L)  at (0,0.5)  {$L$};
  \node[wvtx] (Lp) at (1.55,0.5) {$L^{*}$};
  \node[wvtx] (A)  at (3.0,0.5) {$A$};
  \node[wvtx] (Y)  at (4.4,0.5) {$Y$};
  \draw[ed] (U1) -- (L);
  \draw[ed] (U1) -- (Y);
  \draw[ed] (U2) -- (L);
  \draw[ed] (U2) -- (Lp);
  \draw[ed] (L) -- (Lp);
  \draw[ed] (Lp) -- (A);
  \draw[ed] (A) -- (Y);
\end{tikzpicture}
\captionof{figure}{拥有真实 $L$ 不足以校正混杂，拥有测量到的 $L^{*}$ 却足以校正。}
\label{fig:9.12}
\end{center}
\end{finepoint}

\section{不含测量有误变量的因果图？}

前几章画因果图时，我们一直隐含地作了两条简化假定，二者彼此相关。本节与下一节要把这两条假定及其含义明确说出来。

第一条假定是：图上所有变量都被完美测量。这一假定并不现实，因为在实践中，处理、结局、混杂因子以及其他任何感兴趣的变量，其测量误差往往不可避免。本章描述了如何用因果图表示不同类型测量误差下的测量有误变量，也探讨了使用测量有误的变量（研究者所能得到的也只有这些）识别因果效应的后果。

例如，设我们关心处理 $A$ 对结局 $Y$ 的效应，却只有测量到的处理 $A^{*}$ 与测量到的结局 $Y^{*}$ 的数据。我们已经看到，即使处理 $A$ 对结局 $Y$ 的效应为零，即使既无混杂也无选择偏倚，处理或结局的测量误差也可能在测量到的处理 $A^{*}$ 与测量到的结局 $Y^{*}$ 之间诱发非因果关联。我们也看到，在存在混杂时，混杂因子的测量误差可能使测量到的混杂因子 $L^{*}$ 无法阻断某条后门路径，而若能拿到真实混杂因子 $L$，这条路径本可被成功阻断。

由于可以预期所有变量都会被测量出一些误差，或许有人主张：因果图应始终同时表示其所有变量的真值与测量值。然而在许多情形下，测量误差的大小可以被判定（或已知）小到无关紧要。在这些情形下，一张不区分变量真值与测量值的因果 DAG，也许因其简洁而更可取。如果在所有变量都被完美测量时存在混杂与选择偏倚，那么这些偏倚在存在测量误差时通常也依然存在（例外见精要 9.2）。画不含测量误差的因果图，使我们能专注于混杂与选择偏倚，而不被测量问题分心。

考虑不含测量误差的因果图，往往是很好的第一步近似。一旦对完美测量下可能出现的偏倚有了良好理解，就可以把测量误差作为又一层复杂性加上去。这种两步走的画图办法，有助于我们把两类偏倚——混杂与选择——单独抽出来研究，而不被第三类偏倚——测量——压得喘不过气。本书即循此办法：当重点在混杂与选择偏倚时，我们略去因果图上变量真值与测量值之分。

我们迄今所作的第二条假定同样与测量有关，它是任何因果图的一项基本假定。下一节讨论这一假定。

\section{许多已提出的因果图含有非因果箭头}

考虑图 9.13 的因果图（它等同于图 7.1）。设 $A$、$Y$、$L$ 是三个二值变量，分别代表给 COVID-19 病人的抗病毒治疗、死亡与肥胖（定义为体质指数大于 30）。肥胖病人更可能接受治疗，在未接受治疗时也更可能住院。因此专家画出的 DAG 从 $L$ 出发，分别指向 $A$ 与 $Y$。为简单起见，我们假定是否给治疗只受 $L$ 与医生的偏好影响，而医生偏好与图上其他变量无关。同样为简单起见，假定任何变量都无测量误差。具体地说，$L$ 的真值等于其测量值 $L^{*}$，因此后者无需画在图上。

\begin{figure}[htbp]
\centering
\begin{tikzpicture}[causal,scale=0.95,every node/.style={transform shape}]
  \node[wvtx] (L) at (0,0)    {$L$};
  \node[wvtx] (A) at (1.9,0)  {$A$};
  \node[wvtx] (Y) at (3.7,0)  {$Y$};
  \draw[ed] (L) -- (A);
  \draw[ed] (A) -- (Y);
  \draw[ed] (L.east) .. controls (1.0,1.45) and (2.6,1.45) .. (Y.west);
\end{tikzpicture}
\caption{含非因果箭头 $L\rightarrow Y$ 的“非因果”DAG。}
\label{fig:9.13}
\end{figure}

前文已讨论过，在因果图中处理节点与其他节点地位不同（见 6.4 节）。原因在于，要对处理 $A$ 作用于结局 $Y$ 的效应作有意义的定量因果推断，需要定义良好的反事实结局 $Y^{a}$，而这又要求对处理变量 $A$ 有定义良好的干预（见第 3 章）。在我们的例子里，处理 $A$ 所代表的干预是定义良好的，因为我们清楚抗病毒治疗可以怎样给药或不予给药。因此，从 $A$ 到 $Y$ 的箭头存在与否也是定义良好的。

现在把讨论扩展到对应于非干预的节点，例如图 9.13 中的肥胖 $L$。如第 3 章所述，除非存在定义良好的 ENDE 干预，能在不直接影响结局 $Y$ 的前提下改变非干预变量 $L$，反事实结局 $Y^{l}$ 就是定义不良的。由于 $L$ 是肥胖、$Y$ 是死亡时，多数专家不愿承认存在这样的干预，从 $L$ 到 $Y$ 的箭头就不能被认为定义良好。

本书迄今忽略了这一问题，但健康科学与社会科学中提出的许多 DAG，其箭头的出发变量并不存在定义良好的干预，如图 9.13 中的 $L$。这些箭头，如图 9.13 中的箭头 $L\rightarrow Y$，没有因果解释。为探讨使用这类含非因果箭头的 DAG 的后果，我们现在更细致地讨论图 9.13。我们把箭头 $L\rightarrow A$ 与 $L\rightarrow Y$ 分开来看。

专家画出从 $L\rightarrow A$ 的箭头时，用的是他们的专业知识：他们知道医生一旦得知病人肥胖，就更可能给予治疗。可以设想对 $L$ 作用于 $A$ 的效应存在定义良好的 ENDE 干预，例如给主治医生看一位体质指数与正在作治疗决定的那位病人不同的病人。因此，画从 $L\rightarrow A$ 的箭头是合理的。关于这一点的进一步讨论，见读完下一段之后的精要 9.3。

专家画出从 $L\rightarrow Y$ 的箭头时，用的同样是专业知识：他们知道肥胖病人更可能死亡。然而，由于他们不能接受对 $L$ 作用于 $Y$ 的效应存在 ENDE 干预，箭头 $L\rightarrow Y$ 对他们而言没有因果解释。在本章以及全书各处，我们把“因果 DAG”一词限于所有箭头都有因果解释的 DAG。因此，在目前的知识状态下，图 9.13 是一张“非因果”图。我们把“非因果”加上引号，是为了提醒：因果推断文献中的许多论文仍把同时含因果箭头与非因果箭头的 DAG 定义为因果图。关于“非因果”图的更多讨论见精要 9.4。

\begin{finepoint}[干预是否定义良好取决于所关心的结局]
正文中说，即使对 $L$ 作用于 $Y$ 的效应不存在定义良好的干预，对 $L$ 作用于 $A$ 的效应仍可能存在定义良好的干预。这一说法需要更好的解释：如果我们真的知道怎样对 $L$ 作干预以研究其对 $A$ 的效应，那么有什么能阻止我们用同一干预去研究其对 $Y$ 的效应呢？

这一表面上的矛盾源于记号的不严谨。在我们的例子里，我们用符号 $L$ 指称两个概念：(i) 物理量体重（以公斤计），(ii) 医生所感知到的该量取值（以公斤计）。在 (i) 与 (ii) 数值相同的世界里，混用这两个概念没有实际影响。然而，正文描述对 $L$ 作用于 $A$ 的定义良好的干预时，指的是对 (ii) 而非对 (i) 的干预，因此有必要区分这两个概念。

让我们用符号 $L$ 表示“体重”这一物理量，用符号 $L^{*}$ 表示对 $L$ 的感知。若把 $L^{*}$ 加入因果 DAG，就会有一条从 $L$ 到 $L^{*}$ 的确定性箭头（假定对体重的感知完美），以及一条从 $L^{*}$ 到 $A$ 的常规箭头。按这张扩展后的因果 DAG，体重 $L$ 只有在医生得知取值 $L^{*}$ 时才影响处理 $A$。因此，如果我们对取值 $L^{*}$ 作干预，即使物理量 $L$ 不变，医生的行为也会与我们能以某种方式对 $L$ 作干预时相同。正是在这个意义上，正文说：当结局是 $A$ 时，存在对 $L$ 的定义良好的干预——反事实 $A^{l}$ 定义良好——而当结局是 $Y$ 时则不然——反事实 $Y^{l}$ 定义不良。
\end{finepoint}

如果专家希望画出与“肥胖病人更可能死亡”这一已知事实相容的因果 DAG，他们可以提出图 9.14 的因果图（结构上与图 7.2 相同），其中包含一个对 $L$ 与 $Y$ 都有因果效应的隐藏因素 $H$。节点 $H$ 可以包含未测量的、也许尚未被认识到的变量，例如遗传因素、与体脂储备有关的代谢因素、微生物群，以及大概还有科学尚未发现的其他因素。按图 9.14，$H$ 对处理 $A$ 的效应完全通过 $L$ 中介。也就是说，不存在从 $H$ 直达 $A$ 的箭头。如果例如现实世界作治疗决定所用的唯一信息就是肥胖 $L$，这一假定是合理的。

\begin{figure}[htbp]
\centering
\begin{tikzpicture}[causal,scale=0.95,every node/.style={transform shape}]
  \node[wlat] (H) at (0,-1.7) {$H$};
  \node[wvtx] (L) at (0,0)    {$L$};
  \node[wvtx] (A) at (1.9,0)  {$A$};
  \node[wvtx] (Y) at (3.7,0)  {$Y$};
  \draw[ed] (H) -- (L);
  \draw[ed] (L) -- (A);
  \draw[ed] (A) -- (Y);
  \draw[ed] (H) -- (Y);
\end{tikzpicture}
\caption{加入隐藏因素 $H$ 之后的因果 DAG：$H$ 对 $L$ 与 $Y$ 都有因果效应。}
\label{fig:9.14}
\end{figure}

现在的问题是：图 9.14 中箭头 $H\rightarrow L$ 与 $H\rightarrow Y$ 是否有因果解释，或者说，反事实 $L^{h}$ 与 $Y^{h}$ 是否定义良好。如果专家愿意假定 a) $H$ 对应于一种定义良好的干预，或 b) 存在一种对 $H$ 的定义良好的 ENDE 干预，它只通过影响 $H$ 来影响 $L$ 与 $Y$，那么箭头 $H\rightarrow L$ 与 $H\rightarrow Y$ 就有因果解释。尽管我们对确切的干预 $H$（或对 $H$ 的干预）一无所知，现有知识并不排除其存在，因此箭头 $H\rightarrow L$ 与 $H\rightarrow Y$ 暂时是站得住脚的，我们也就把图 9.14 视为因果图。

\begin{finepoint}[具有定义良好的统计解释的“非因果”图]
考虑某个 FFR-CISTG 模型的 DAG 表示，它假定 (i) 处理 $A$ 是唯一具有定义良好的干预、从而有反事实 $(M^{a},Y^{a})$ 的变量；(ii) 图上变量的分布按该 DAG 分解，即每个变量给定其父节点后与其非后代条件独立。在这类 DAG 中，箭头一般没有因果解释。箭头只是通过 d-分离，编码图上变量以及相应 SWIG 所满足的条件独立（见技术注记 21.12）。在这种非因果的 DAG 解释下，图 9.13 中的箭头 $L\rightarrow Y$ 不必仅因对 $L$ 的干预定义不良就被删去。

Richardson and Robins (2013) 指出这一进路有严重困难。具体地说，如果箭头没有因果解释，就没有理由指望研究变量的分布可以由任何（不完全的）DAG 表示（即按该 DAG 分解）。存在未测变量 $U$ 时，这一困难会放大，因为那时甚至无法从观测数据中经验地检验某些条件独立。例如，图 7.14 的前门图蕴含：给定 $M$ 与 $U$ 时 $Y$ 与 $A$ 独立。任何选择图 7.14 作为合适因果 DAG 的研究者，必定有实质性的理由来设定这一条件独立。的确，许多研究者很难想象，除了相信 $A$ 与 $Y$ 共有原因 $U$、且 $A$ 对 $Y$ 的因果效应完全由 $M$ 中介之外，还能有什么别的理由——而这一信念恰恰赋予图上的每条箭头以因果解释。另见精要 9.5。

另一种看法是，把非因果箭头的使用视为研究者对“$M$-反事实 $Y^{m}$ 存在”这一强因果主张持怀疑态度时的一种回应。这种“非因果”进路把图 7.14 解释为表示观测变量之间所假定的统计独立 $Y\indep A\given M$，而这一独立将在未来一项对 $A$ 作随机分配的试验中成立。因此，如果这一独立在未来试验中不成立，那么“$M$-反事实存在”这一强因果主张的依据，就与图 7.14 的结构一起被证伪了。
\end{finepoint}

\section{许多已提出的图含有非因果箭头，这要紧吗？}

我们已经看到，专家最初提出的 DAG（图 9.13）不是因果 DAG，因为其中一条箭头（从 $L$ 到 $Y$ 的那条）无法作因果解释。然而，尽管在因果上是错的，这张 DAG 对于关于处理 $A$ 作用于结局 $Y$ 之效应的因果推断却是够用的：无论我们用图 9.13 的非因果 DAG，还是用图 9.14 的因果 DAG，都得出结论——$A$ 与 $Y$ 之间的所有后门路径都通过对 $L$ 条件化而阻断。

因此，图 9.13 的 DAG 虽然缺少节点 $U$ 并含有非因果箭头 $L\rightarrow Y$，却正确地指导了数据分析，因为对我们的例子中 $A$ 对 $Y$ 的平均因果效应而言，校正 $L$ 正是识别公式——标准化或逆概率加权——所需的全部。这一结论在预料之中，因为两张 DAG 中都没有从不可测变量指向处理 $A$ 的箭头（另见精要 9.2）。这一过度简化的情景说明了一个普遍问题：在现实的复杂情形中，专家知识很少好到足以画出一张各组成部分全都已知的因果图。

事实上，健康科学与社会科学中提出的许多 DAG 其实是非因果 DAG，因为它们缺少使 DAG 成为因果图的隐藏变量 $H$。把一个对其后代的效应没有定义良好干预的变量（如肥胖及其后代死亡）纳入图中，我们实际上是在宣告自己的 DAG 是非因果的。不过，也有可能：从非因果 DAG 与从含隐藏节点的因果 DAG 推导出的、关于所关心因果效应的识别公式相同。我们的例子正是如此：如果图 9.14 才是正确的 DAG，那么使用图 9.13 的 DAG 会得到相同的识别公式。

但非因果 DAG 可能误导人。考虑图 9.15 的 DAG。宣告这张 DAG 是因果 DAG，意味着我们相信对 $L$ 作用于 $Y$ 的效应存在定义良好的干预（因为有一条从 $L$ 直达 $Y$ 的箭头）。反之，若不存在这样的干预，该 DAG 就是非因果的，我们只能指望基于 $L$ 的识别公式恰好正确。

\begin{figure}[htbp]
\centering
\begin{tikzpicture}[causal,scale=0.95,every node/.style={transform shape}]
  \node[wlat] (U) at (0,-1.5)  {$U$};
  \node[wvtx] (L) at (0.15,0.1) {$L$};
  \node[wvtx] (A) at (2.1,0)   {$A$};
  \node[wvtx] (Y) at (4.0,0)   {$Y$};
  \draw[ed] (U) -- (L);
  \draw[ed] (U) -- (A);
  \draw[ed] (A) -- (Y);
  \draw[ed] (L.east) .. controls (1.2,1.5) and (2.9,1.5) .. (Y.west);
\end{tikzpicture}
\caption{若宣告为因果 DAG，则须相信对 $L$ 作用于 $Y$ 存在定义良好的干预。}
\label{fig:9.15}
\end{figure}

循上述推理，图 9.16 描绘的情形是：我们测量到的变量 $L$ 其实是一个未知变量 $H$ 的替代，而对 $H$ 存在定义良好的干预。在图 9.15 中，$A$ 与 $Y$ 之间的后门路径被 $L$ 阻断，而在图 9.16 中则没有。因此，如果研究者不知道自己所测量的变量只是替代混杂因子 $L$，从而提出图 9.15 而非图 9.16，他们就会得出错误结论，以为 $L$ 足以阻断所有后门路径——如 9.3 节所讨论。画因果 DAG 时，我们需要仔细思量：碰巧被测量的变量，是否也就是存在定义良好干预的那些变量。否则，对 DAG 中非因果箭头的疏忽，会让我们对效应估计的有效性产生虚假的信心。另一个例子见精要 9.6。

\begin{figure}[htbp]
\centering
\begin{tikzpicture}[causal,scale=0.92,every node/.style={transform shape}]
  \node[wlat] (U) at (0,-2.0)  {$U$};
  \node[wlat] (H) at (0.4,0.3) {$H$};
  \node[wvtx] (L) at (0.4,1.8) {$L$};
  \node[wvtx] (A) at (2.5,0.3) {$A$};
  \node[wvtx] (Y) at (4.4,0.3) {$Y$};
  \draw[ed] (U) -- (H);
  \draw[ed] (U) -- (A);
  \draw[ed] (H) -- (L);
  \draw[ed] (A) -- (Y);
  \draw[ed] (H.east) .. controls (1.6,1.75) and (3.4,1.75) .. (Y.west);
\end{tikzpicture}
\caption{测量到的 $L$ 只是未知变量 $H$ 的替代。}
\label{fig:9.16}
\end{figure}

自第 6 章引入因果图以来，我们已经走了很远。因果 DAG 与 SWIG 是研究者组织和交流其因果假定的强大工具，但如本章所见，这些图也受制于因果推断普遍固有的那些实际约束。具体地说，若对 $X$ 作用于 $Y$ 的效应没有定义良好的干预，因果箭头 $X\rightarrow Y$ 就无法被有意义地解释。

提出一张因果 DAG 时，我们需要仔细思考其中每条箭头的解释。对含非因果箭头的 DAG 作科学上不加审视的接受，可能导致因果推断的错误结论。对于表示电路（其中所有干预都定义良好）的因果图，这种审视没有必要；但对于健康科学与社会科学中提出的因果图，则不可或缺。

为教学目的，上述讨论简化了定义良好的干预这一概念。如第 3 章所述，没有哪种干预是完美定义良好的，但对某些干预，科学界的共识是它们已足够定义良好。我们已论证，定量因果推断从根本上依赖于“足够定义良好的干预”这一（诚然模糊的）概念。因此，在本书余下部分，除少数明确标示的例外，我们将假定所有 DAG 都是严格因果的，即 DAG 上的每条箭头都代表一种足够定义良好的干预。也就是说，每条箭头都联系着一种可以如此规定的干预：依现有科学知识，不再留有有意义的含糊。几年或几十年后，我们也许会发现自己原先的信念、从而我们的因果图，是错的。

\begin{finepoint}[与前门公式的联系]
图 9.14 正是图 7.14 的前门图，只是把 $L\rightarrow A$ 替换成了 $A\rightarrow M$。因此，$\E[Y^{l}]$ 可由技术注记 7.4 的前门公式识别，只需把 $L$、$l$、$l'$ 替换为 $A$、$a$、$a'$，把 $A$、$a$ 替换为 $M$、$m$。

设一位研究者修改图 9.14，加上一条直达的 $L\rightarrow Y$ 箭头。理由在于：当营养师与理疗师接诊某位 $L$ 取值很高的个体时，会提供额外的辅助照护（如饮食建议、预防深静脉血栓的运动等），而这些并未记录在研究者可获得的病历中。未修改的图 9.14 中的反事实变量 $Y^{l}$，与加上直达 $L\rightarrow Y$ 箭头后修改图中的变量 $Y^{l}$ 不同。在未修改的图中，对比 $Y^{l}-Y^{l'}$ 是 $l$ 相对于 $l'$ 经由因果通路 $L\rightarrow A\rightarrow Y$ 的因果效应。相反，在修改后的图中，$Y^{l}-Y^{l'}$ 是沿 $L\rightarrow Y$ 与 $L\rightarrow A\rightarrow Y$ 两条通路的总效应。事实上，与未修改的图不同，在修改后的图下 $\E[Y^{l}]$ 不可识别。

人们或许会猜想：$l$ 相对于 $l'$ 沿通路 $L\rightarrow A\rightarrow Y$ 的效应，在两张图中应当相同，因而仍可由前门公式识别。这一猜想确实成立，但证明要留到第 23 章研究通路特定效应（path-specific effects）的识别时再给出。
\end{finepoint}

\begin{center}
\begin{tikzpicture}[causal,scale=0.68,every node/.style={transform shape}]
  \node[wlat] (U) at (0,-1.6)  {$U$};
  \node[wvtx] (A) at (0.2,0)   {$A$};
  \node[wvtx] (L) at (2.2,0)   {$L$};
  \node[wvtx] (Y) at (4.1,0)   {$Y$};
  \draw[ed] (U) -- (A);
  \draw[ed] (U) -- (L);
  \draw[ed] (A) -- (L);
  \draw[ed] (L) -- (Y);
  \draw[ed] (A.east) .. controls (1.3,1.5) and (3.0,1.5) .. (Y.west);
\end{tikzpicture}
\captionof{figure}{研究者最初提出的 DAG：$L$ 是 $A$ 的后代。}
\label{fig:9.17}
\vspace{10pt}
\begin{tikzpicture}[causal,scale=0.68,every node/.style={transform shape}]
  \node[wlat] (U) at (0,-1.7)  {$U$};
  \node[wvtx] (A) at (0.2,0)   {$A$};
  \node[wlat] (H) at (2.2,0)   {$H$};
  \node[wvtx] (L) at (2.2,-1.7){$L$};
  \node[wvtx] (Y) at (4.1,0)   {$Y$};
  \draw[ed] (U) -- (A);
  \draw[ed] (U) -- (H);
  \draw[ed] (A) -- (H);
  \draw[ed] (H) -- (Y);
  \draw[ed] (H) -- (L);
  \draw[ed] (A.east) .. controls (1.3,1.5) and (3.0,1.5) .. (Y.west);
\end{tikzpicture}
\captionof{figure}{加入隐藏变量 $H$、把 $L$ 表示为 $H$ 的替代之后的 DAG。}
\label{fig:9.18}
\vspace{10pt}
\begin{tikzpicture}[causal,scale=0.68,every node/.style={transform shape}]
  \node[wlat] (U) at (0,-1.7)  {$U$};
  \node[wvtx] (A) at (0.2,0)   {$A$};
  \node[wlat] (H) at (2.2,0)   {$H$};
  \node[wvtx] (L) at (2.2,-1.7){$L$};
  \node[wvtx] (Y) at (4.1,0)   {$Y$};
  \draw[ed] (U) -- (A);
  \draw[ed] (U) -- (L);
  \draw[ed] (A) -- (H);
  \draw[ed] (H) -- (Y);
  \draw[ed] (H) -- (L);
  \draw[ed] (A.east) .. controls (1.3,1.5) and (3.0,1.5) .. (Y.west);
\end{tikzpicture}
\captionof{figure}{$U$ 直接影响替代变量 $L$ 的因果 DAG。}
\label{fig:9.19}
\end{center}

\begin{finepoint}[从非因果图到因果图]
设某些关心 $A$ 对 $Y$ 之效应的研究者提出了图 9.17 的 DAG。画这张 DAG 时，他们依据的是关于变量时间先后顺序的先验知识，以及以下两个事实：(i) 测量到的变量 $L$ 与 $Y$ 有关联；(ii) 存在一组未测量但已知的因素 $U$，它影响 $A$，并与 $L$ 有关联。如精要 7.4 所述，若图 9.17 是真实的因果图，则 $A$ 对 $Y$ 的效应不可识别，因为 $L$ 是 $A$ 的后代。

进一步思索后，研究者意识到对 $L$ 作用于 $Y$ 的效应不存在定义良好的干预。因此箭头 $L\rightarrow Y$ 不是因果箭头，他们的 DAG 不是因果的。为把图 9.17 变成因果图，他们加入隐藏变量 $H$，画一条指向 $Y$ 的因果箭头，并把 $L$ 表示成 $H$ 的替代。图 9.18 描绘了他们修改后的 DAG。若图 9.18 是真实的因果图，则 $A$ 对 $Y$ 的效应不可识别，因为后门路径 $A\leftarrow U\rightarrow H\rightarrow Y$ 无法被任何测量到的变量阻断。

注意，研究者保留了箭头 $U\rightarrow A$，这意味着他们相信对 $U$ 作用于 $A$ 的效应存在定义良好的干预。他们又把图 9.17 中从 $U$ 与 $A$ 指向 $L$ 的箭头，改为在图 9.18 中指向 $H$。这意味着研究者相信：对 $U$ 与 $A$ 作用于 $H$ 的效应必定存在定义良好的干预，且 $U$ 与 $A$ 对 $L$ 的效应完全由它们对 $H$ 的效应中介。（即使 $L$ 是 $H$ 的确定性函数——这与图 9.18 相容——对 $L$ 条件化也不会阻断经过 $H$ 的路径，因为 $H$ 不是 $L$ 的确定性函数：$H$ 的维数与复杂性都高于 $L$。）

图 9.18 中 $H$ 继承图 9.17 中所有指向 $L$ 的箭头，这并不总是站得住脚。例如，$U$ 可能像图 9.19 的因果 DAG 那样直接影响 $L$，而不是像图 9.18 那样影响 $H$。事实上，若图 9.19 才是真实的因果 DAG，则 $A$ 对 $Y$ 的效应可识别，因为 $A$ 与 $Y$ 之间没有开放的后门路径。注意，在图 9.19 中，替代变量 $L$ 不可能是 $H$ 的确定性函数，因为它还受 $U$ 影响。与图 9.19 一样、$U$ 对替代变量 $L$ 有直接效应而对 $H$ 无直接效应的一个例子如下：$U$ 表示医生决定开具某项诊断检查，$L$ 表示检查结果，$H$ 表示该检查结果背后的生物学决定因素。
\end{finepoint}
